\chapter{Testovanie}

\section{Heuristická analýza}


\section{Používateľské testovanie}
Používateľské testovanie som realizoval v troch fázach. Prvou bolo testovanie Lo-Fi prototypu, a následne dve fázy testovania Hi-Fi prototypu. 

\subsection{Dotazník pred testovaním}
Na zaradenie testovacích subjektov do jednotlivých vekových kategórií a typu používateľa som použil dotazník, ktorý obsahoval otázky týkajúce sa veku, skúseností s technológiami, záujmu o skauting a aktívnu participáciu v rámci skautingu. Tento dotazník mi pomohol získať základné informácie o testovacích subjektoch a zabezpečiť, že testovanie bude reprezentatívne pre cieľovú skupinu aplikácie.

\begin{enumerate}
    \item Meno a priezvisko
    \item Vek
    \item Skúsenosti s technológiami (počet rokov používania smartfónu, používanie aplikácií, atď.)
    \item Záujem o skauting (člen skautskej organizácie)
    \item Aktívna participácia v rámci skautingu (plnenie programu, radcovanie)
\end{enumerate}

\subsection{Dotazník po testovaní}
Na získanie spätnej väzby od testovacích subjektov som použil dotazník, ktorý obsahoval otázky týkajúce sa používateľského zážitku, navigácie a celkového dojmu z aplikácie. 

\begin{enumerate}
    \item Ako by ste ohodnotili celkový zážitok z používania aplikace?
    \item Bolo pre vás jednoduché navigovat mezi jednotlivými obrazovkami?
    \item Aké prvky používateľského rozhrania se vám páčili a ktoré nie?
    \item Máte nejaké návrhy na zlepšenie aplikácie?
\end{enumerate}

\subsection{Testovanie Lo-Fi prototypu}
Testovanie prebiehalo s jedným testovacím subjektom a zameriavalo sa na logické členenie aplikácie a navigáciu. Testovací subjekt bol požiadaný, aby vykonal sériu úloh, ktoré zahrňovali navigáciu medzi rôznymi obrazovkami aplikácie a interakciu s jednotlivými prvkami používateľského rozhrania.

\subsection{Testovanie Hi-Fi prototypu}
Testovanie prebiehalo v dvoch fázach, pričom v každej fáze sa zúčastnilo 5 testovacích subjektov. Testovacie subjekty boli požiadané, aby vykonali sériu úloh, ktoré zahrňovali bežné používanie aplikácie, ako je prehliadanie obsahu, plnenie úloh, pridávanie sa do družiny a sledovanie pokroku členov družiny. Počas testovania som pozoroval interakciu testovacích subjektov s aplikáciou a zaznamenával som ich spätnú väzbu, ktorú som následne analyzoval a použil na vylepšenie aplikácie.

\subsection{Testovací scénar}


Testovací scénar pre radcov:
\begin{enumerate}
    \item Vytvorte si účet a založte novú družinu. 
    \item Zobrazte kód družiny na pridanie členov.
    \item Zobrazte zoznam členov družiny.
    \item Na schôdzke ste s všetkými členmi sružiny splnili úlohu Identita z Prvého skauta. Splňte túto úlohu členom družiny v aplikácii.
    \item Blíži sa koniec roka a chcete zistiť komu sa podarilo splniť úlohu Tik-tak z Prvého skauta. Zobrazte progres všetkých členov družiny v tejto úlohe.
    \item Člen vašej družiny Breza nemá smartphone, ale chcete evidovat jeho progres. Vytvorte virtuálneho člena družiny s prezývkou Breza. (Pridané do druhej fázy testovania)
\end{enumerate}

Testovací scénar pro členov:
\begin{enumerate}
    \item Vytvorte si účet a přidejte se do družiny pomocou kódu družiny.
    \item Chcete zistiť aké úlohy sú potrebné na získanie zeleného stupňa odborky táborník. Pridajte si ich medzi sledované.
    \item Označte úlohy na splnenie zeleného stupňa odborky táborník za splnené.
    \item Označte úlohu Spoza hraníc z Prvého skauta za splnenú a požiadajte vedúceho družiny o schválenie úlohy.
    \item Popri vašom programe sa chcete venovat aj radcovaniu. Založte si vlastnú družinu.
\end{enumerate}

\section{Prvá fáza používateľského testovania}
Prvá fáza používateľského testovania sa zúčastnilo 6 subjektov naprieč vekovým rozsahom 11 až 18 rokov. Testovanie sa uskutočnilo v priestoroch zborovej klubovne skautského zboru 96. Malokarpatský zbor - Bratislava Rača nakoľko sa jednalo o prostredie v ktorom sa subjekty bežne zdržiavali pri skautských schôdzkach. Pred testovaním boli testeri oboznámeny s ideou testavania a so skutočnosťou že zlyhanie pri plnení scénara nie je ich chyba ale chyba aplikácie.

\subsection*{Tester 1: Kristián}
\subsubsection*{Dotazník pred testovaním}
\begin{itemize}[leftmargin=*]
    \item \textbf{Vek:} 12
    \item \textbf{Skúsenosti s technologiámi (počet rokov používania smartphonu):} 6
    \item \textbf{Záujem o skauting (člen skautskej organizácie):} Pravidelná účasť na zborových podujatiach
    \item \textbf{Aktívna participácia v rámci skautingu:} Program, začiatky radcovania
\end{itemize}

\subsubsection*{Splnenie scénarov}
\begin{table}[h!]
    \centering
    \begin{tabular}{|l|p{1.5cm}|p{1.5cm}|p{1.5cm}|p{1.5cm}|p{1.5cm}|}
        \hline
        Scénar & 1 & 2 & 3 & 4 & 5 \\ \hline
        Splnenie & & & & & \\ \hline
    \end{tabular}
    \caption{Tester 1 - splnenie scénarov}
\end{table}

\subsubsection*{Dotazník po testovaní}
\begin{enumerate}[leftmargin=*]
    \item \textbf{Ako by ste ohodnotili celkový zážitok z používania aplikácie?}
    \begin{itemize}
        \item 10
    \end{itemize}
    \item \textbf{Bolo pre vás jednoduché navigovať medzi jednotlivými obrazovkami?}
    \begin{itemize}
        \item stredne
    \end{itemize}
    \item \textbf{Aké prvky používateľského rozhrania se vám páčili a ktoré nie?} kategórie
    \item \textbf{Máte nejaké návrhy na zlepšenie aplikácie?} splnene ulohy by si mal rovno roztiahnute
\end{enumerate}

\subsection*{Tester 2: Alex}
\subsubsection*{Dotazník pred testovaním}
\begin{itemize}[leftmargin=*]
    \item \textbf{Vek:} 13
    \item \textbf{Skúsenosti s technológiami (počet rokov používania smartfónu):} 4
    \item \textbf{Záujem o skauting (člen skautskej organizácie):} Pravidelná účasť na zborových podujatiach
    \item \textbf{Aktívna participácia v rámci skautingu:} Program
\end{itemize}

\subsubsection*{Splnenie scénarov}
\begin{table}[h!]
    \centering
    \begin{tabular}{|l|p{1.5cm}|p{1.5cm}|p{1.5cm}|p{1.5cm}|p{1.5cm}|}
        \hline
        Scénar & 1 & 2 & 3 & 4 & 5 \\ \hline
        Splnenie & & & & & \\ \hline
    \end{tabular}
    \caption{Tester 2 - splnenie scénarov}
\end{table}

\subsubsection*{Dotazník po testovaní}
\begin{enumerate}[leftmargin=*]
    \item \textbf{Ako by ste ohodnotili celkový zážitok z používania aplikácie?}
    \begin{itemize}
        \item 8
    \end{itemize}
    \item \textbf{Bolo pre vás jednoduché navigovať medzi jednotlivými obrazovkami?}
    \begin{itemize}
        \item ano
    \end{itemize}
    \item \textbf{Aké prvky používateľského rozhrania sa vám páčili a ktoré nie?} odborky a ulohy
    \item \textbf{Máte nejaké návrhy na zlepšenie aplikácie?} tlacitka nie su jasne (or podobne)
\end{enumerate}

\subsection*{Tester 3: Tadeáš}
\subsubsection*{Dotazník pred testovaním}
\begin{itemize}[leftmargin=*]
    \item \textbf{Vek:} 18
    \item \textbf{Skúsenosti s technológiami (počet rokov používania smartfónu):} 9
    \item \textbf{Záujem o skauting (člen skautskej organizácie):} Pravidelná účasť na zborových podujatiach
    \item \textbf{Aktívna participácia v rámci skautingu:} Program, radcovanie
\end{itemize}

\subsubsection*{Splnenie scénarov}
\begin{table}[h!]
    \centering
    \begin{tabular}{|l|p{1.5cm}|p{1.5cm}|p{1.5cm}|p{1.5cm}|p{1.5cm}|}
        \hline
        Scénar & 1 & 2 & 3 & 4 & 5 \\ \hline
        Splnenie & & & & & \\ \hline
    \end{tabular}
    \caption{Tester 3 - splnenie scénarov}
\end{table}

\subsubsection*{Dotazník po testovaní}
\begin{enumerate}[leftmargin=*]
    \item \textbf{Ako by ste ohodnotili celkový zážitok z používania aplikácie?}
    \begin{itemize}
        \item 7
    \end{itemize}
    \item \textbf{Bolo pre vás jednoduché navigovať medzi jednotlivými obrazovkami?}
    \begin{itemize}
        \item nie
    \end{itemize}
    \item \textbf{Aké prvky používateľského rozhrania sa vám páčili a ktoré nie?} vyzvy a odborky
    \item \textbf{Máte nejaké návrhy na zlepšenie aplikácie?} tutorial
\end{enumerate}

\subsection*{Tester 4: Fabián}
\subsubsection*{Dotazník pred testovaním}
\begin{itemize}[leftmargin=*]
    \item \textbf{Vek:} 17
    \item \textbf{Skúsenosti s technológiami (počet rokov používania smartfónu):} 9
    \item \textbf{Záujem o skauting (člen skautskej organizácie):} Pravidelná účasť na zborových podujatiach
    \item \textbf{Aktívna participácia v rámci skautingu:} Program, radcovanie
\end{itemize}

\subsubsection*{Splnenie scénarov}
\begin{table}[h!]
    \centering
    \begin{tabular}{|l|p{1.5cm}|p{1.5cm}|p{1.5cm}|p{1.5cm}|p{1.5cm}|}
        \hline
        Scénar & 1 & 2 & 3 & 4 & 5 \\ \hline
        Splnenie & & & & & \\ \hline
    \end{tabular}
    \caption{Tester 4 - splnenie scénarov}
\end{table}

\subsubsection*{Dotazník po testovaní}
\begin{enumerate}[leftmargin=*]
    \item \textbf{Ako by ste ohodnotili celkový zážitok z používania aplikácie?}
    \begin{itemize}
        \item 8
    \end{itemize}
    \item \textbf{Bolo pre vás jednoduché navigovať medzi jednotlivými obrazovkami?}
    \begin{itemize}
        \item ano
    \end{itemize}
    \item \textbf{Aké prvky používateľského rozhrania sa vám páčili a ktoré nie?} prezeranie profilu
    \item \textbf{Máte nejaké návrhy na zlepšenie aplikácie?} zlepsit co su tlacitka
\end{enumerate}

\subsection*{Tester 5: Filip}
\subsubsection*{Dotazník pred testovaním}
\begin{itemize}[leftmargin=*]
    \item \textbf{Vek:} 16
    \item \textbf{Skúsenosti s technológiami (počet rokov používania smartfónu):} 5
    \item \textbf{Záujem o skauting (člen skautskej organizácie):} Pravidelná účasť na zborových podujatiach
    \item \textbf{Aktívna participácia v rámci skautingu:} Program, radcovanie
\end{itemize}

\subsubsection*{Splnenie scénarov}
\begin{table}[h!]
    \centering
    \begin{tabular}{|l|p{1.5cm}|p{1.5cm}|p{1.5cm}|p{1.5cm}|p{1.5cm}|}
        \hline
        Scénar & 1 & 2 & 3 & 4 & 5 \\ \hline
        Splnenie & & & & & \\ \hline
    \end{tabular}
    \caption{Tester 5 - splnenie scénarov}
\end{table}

\subsubsection*{Dotazník po testovaní}
\begin{enumerate}[leftmargin=*]
    \item \textbf{Ako by ste ohodnotili celkový zážitok z používania aplikácie?}
    \begin{itemize}
        \item 9
    \end{itemize}
    \item \textbf{Bolo pre vás jednoduché navigovať medzi jednotlivými obrazovkami?}
    \begin{itemize}
        \item ano
    \end{itemize}
    \item \textbf{Aké prvky používateľského rozhrania sa vám páčili a ktoré nie?} program je prehladny
    \item \textbf{Máte nejaké návrhy na zlepšenie aplikácie?} kliknutie, potvrdzovanie
\end{enumerate}

\subsection*{Tester 6: Zina}
\subsubsection*{Dotazník pred testovaním}
\begin{itemize}[leftmargin=*]
    \item \textbf{Vek:} 16
    \item \textbf{Skúsenosti s technológiami (počet rokov používania smartfónu):} 7
    \item \textbf{Záujem o skauting (člen skautskej organizácie):} Pravidelná účasť na zborových podujatiach
    \item \textbf{Aktívna participácia v rámci skautingu:} Program, radcovanie
\end{itemize}

\subsubsection*{Splnenie scénarov}
\begin{table}[h!]
    \centering
    \begin{tabular}{|l|p{1.5cm}|p{1.5cm}|p{1.5cm}|p{1.5cm}|p{1.5cm}|}
        \hline
        Scénar & 1 & 2 & 3 & 4 & 5 \\ \hline
        Splnenie & & & & & \\ \hline
    \end{tabular}
    \caption{Tester 6 - splnenie scénarov}
\end{table}

\subsubsection*{Dotazník po testovaní}
\begin{enumerate}[leftmargin=*]
    \item \textbf{Ako by ste ohodnotili celkový zážitok z používania aplikácie?}
    \begin{itemize}
        \item 9
    \end{itemize}
    \item \textbf{Bolo pre vás jednoduché navigovať medzi jednotlivými obrazovkami?}
    \begin{itemize}
        \item stredne
    \end{itemize}
    \item \textbf{Aké prvky používateľského rozhrania sa vám páčili a ktoré nie?} zobrazenie profilu
    \item \textbf{Máte nejaké návrhy na zlepšenie aplikácie?} jasnejsie, tutorial
\end{enumerate}

\subsection{Problémy odhalené počas prvej fáze testovania}
Pri testovaní sa objavovali niektoré chyby pravidelne. Medzi hlavné patrí zlá rozličnosť pri akčných prvkoch. Tlačítka ktoré obsahovali ikony boli neprehľadné a splávali s textom, využitie vertikálneho škálovania tlačítiek by mohol napomôcť k vyriešeniu tohto problému. Ďalšou problematickou oblasťou boli nejasné akčné prvky v podobe gest, pridanie tlačítka by obmedzilo využitie priestoru na obrazovke, lepšou alternatívou by bolo pridanie jednorázového návodu na použitie pri prvom spustení aplikácie.

Po uskutočnení testovania na kontrolovanom zariadení bolo umožnené participantom staihnúť si aplikáciu na vlastné zariadenie. Rôzne veľkosti zariadení participantov identifikovali problém s chýbajúcim scrollovaním na niektorých obrazovkách pri menších veľkostiach displejov. 

Problém s offline režimom.

Vytvorenie virtuálneho člena družiny.


\subsection{Problémy odhalené počas druhej fáze testovania}
Zabudnuté heslo.



\section{Automatizované testovanie}



% \section{Testovacie scenáre pre radcov}

% \subsection{Testovací scénar 1: Vytvorenie účtu a vytvorenie družiny}
% Vytvorte si účet a založte novú družinu. 

% \subsection{Testovací scénar 2: Pridanie členov do družiny}
% Zobrazte kód družiny na pridanie členov.

% \subsection{Testovací scénar 3: Zobrazenie členov družiny}
% Zobrazte zoznam členov družiny.

% \subsection{Testovací scénar 4: Splnenie úlohy členom družiny}
% Na schôdzke ste s všetkými členmi sružiny splnili úlohu Identita z Prvého skauta. Splňte túto úlohu členom družiny v aplikácii.

% \subsection{Testovací scénar 5: Zobrazenie progresu členov družiny}
% Blíži sa koniec roka a chcete zistiť komu sa podarilo splniť úlohu Tik-tak z Prvého skauta. Zobrazte progres všetkých členov družiny v tejto úlohe.

% \section{Testovacie scénare pre členov}

% \subsection{Testovací scénar 1: Vytvorenie účtu a přidanie sa do družiny}
% Vytvorte si účet a přidejte se do družiny pomocou kódu družiny.

% \subsection{Testovací scénar 2: Prezeranie programu a sledovanie úloh}
% Chcete zistiť aké úlohy sú potrebné na získanie zeleného stupňa odborky táborník. Pridajte si ich medzi sledované.

% \subsection{Testovací scénar 3: Splnenie úlohy}
% Označte úlohy na splnenie zeleného stupňa odborky táborník za splnené.

% \subsection{Testovací scénar 4: Splnenie úlohy so schválením vedúceho družiny}
% Označte úlohu Spoza hraníc z Prvého skauta za splnenú a požiadajte vedúceho družiny o schválenie úlohy.

% \subsection{Testovací scénar 5: Začatie radcovania}
% Popri vašom programe sa chcete venovat aj radcovaniu. Založte si vlastnú družinu.